\documentclass[11pt,a4paper]{article}
\usepackage[utf8]{inputenc}
\usepackage{amsmath,amssymb,amsfonts}
\usepackage{booktabs}
\usepackage{longtable}
\usepackage{array}
\usepackage{geometry}
\usepackage{tikz}
\usepackage{hyperref}
\usepackage{float}

\usetikzlibrary{shapes.geometric, arrows, positioning, fit, backgrounds}

\geometry{margin=1in}

\title{\textbf{Model Description of the Elemental World (\texttt{ODD 2020 Protocol})}}
\author{\textbf{Filipe M. Andrade, Roberto Molowny-Horas, Marcos Fernandez-Martinez}}
\date{\today}

\begin{document}

\maketitle

\abstract{

}

\section{Overview}

\subsection{Purpose and Patterns}
The primary purpose of the \texttt{HybridEcosystem} model is to investigate how individual elemental composition  adaptation, multi-elemental resource competition, local soil nutrient cycling, and spatial environmental heterogeneity interact to govern species coexistence, community assembly and ecosystem resilience under disturbance regimes.

The model is designed to observe and reproduce several ecological patterns:

\begin{enumerate}
    \item Species coexistence and niche differentiation based on multi-elemental composition optima.
    \item Spatial self-organization of plant biomass driven by localized competition and seed dispersal constraints.
    \item Stoichiometric limitation and quota-based growth dynamics following Liebig's Law of the Minimum.
    \item Ecosystem collapse, recovery, and compositional shifts under periodic environmental catastrophes and strength-dependent disturbance filters.
\end{enumerate}

\subsection{Entities, State Variables, and Scales}

\subsubsection{Entities}
The model consists of three core entity types:
\begin{enumerate}
    \item \textbf{Individual Agents (Plants):} Mobile/spatial individual organisms buffer holding up to $M = \text{\texttt{MAX\_AGENTS}}$ entries.
    \item \textbf{Grid Cells (Soil Environment):} A 2D regular spatial grid of size $H \times W$ containing inorganic and organic soil nutrient pools.
    \item \textbf{Global Macro-Environment:} Global temperature dynamics and spatial noise fields.
\end{enumerate}

\subsubsection{Spatial and Temporal Scales}
\begin{itemize}
    \item \textbf{Spatial Domain:} A two-dimensional grid of height $H$ cells and width $W$ cells. Toroidal boundary conditions (periodic boundaries) apply during seed dispersal.
    \item \textbf{Spatial Scale:} 1 grid cell represents a localized spatial unit holding plant biomass and soil nutrient pools (spatial cell area is dimensionless/arbitrary).
    \item \textbf{Temporal Domain:} Discrete time steps $t \in \{0, 1, 2, \dots, T\}$. Each time step represents an arbitrary ecological update period (e.g., day).
\end{itemize}

\subsection{Process Overview and Scheduling}
Within every discrete time step $t$, the simulation executes processes in the following synchronous and sequential order:

\begin{enumerate}
    \item \textbf{Soil Inorganic Diffusion and Inflow:} Soil inorganic pools diffuse via a $3 \times 3$ depthwise spatial kernel. Stochastic atmospheric nutrient inflow with potential pulse shocks is added.
    \item \textbf{Environmental Field Updating:} Persistent environmental ($\mathbf{E}$) and shock ($\mathbf{S}$) fields update via autoregressive smoothed Gaussian noise, and global temperature $T(t)$ updates sinusoidally.
    \item \textbf{Localized Patch Disturbance:} Random spatial disturbance events deplete local inorganic soil nutrient pools within a radius $R_{\text{dist}}$.
    \item \textbf{Agent Niche Fitness Calculation:} Each active agent evaluates its elemental niche fitness $w_i(t)$ based on the Mahalanobis distance between its internal elementome ratio $\boldsymbol{q}_i$ and its species niche center $\boldsymbol{c}_{s_i}$.
    \item \textbf{Nutrient Uptake and Limitation:} Agents determine growth potential, take up carbon internally and draw inorganic soil nutrients via Michaelis-Menten kinetics. Growth is constrained by element cell quotas (Liebig's Law) and local crowding/biomass capacity $K_{\text{bio}}$.
    \item \textbf{Biomass Update, Respiration, and Demographic Stochasticity:} Realized biomass growth is added, temperature-dependent respiration is subtracted and demographic noise is applied.
    \item \textbf{Mortality and Disturbance Filters:} Agents undergo viability filters, periodic catastrophes, weak niche-dependent disturbances and strong mass-dependent disturbances.
    \item \textbf{Turnover, Recycling, and Mineralization:} Biomass from dead agents and living turnover is recycled into local organic soil pools. Organic soil pools undergo temperature-dependent mineralization into inorganic soil pools.
    \item \textbf{Reproduction and Dispersal:} Eligible fertile agents spend biomass $C_{\text{seed}}$ to generate seeds, which disperse within a species-specific spatial radius $R_{s_i}$ and establish based on local density constraints.
\end{enumerate}

\section{Design Concepts}

\subsection{Basic Principles}
The model synthesizes fundamental principles from elemental ecology, niche theory, and individual-based ecological dynamics. Growth is governed by internal elemental ratios relative to species-specific niche optima. Resource uptake follows Michaelis-Menten saturation kinetics, while growth bottlenecks obey Liebig's Law of the Minimum. Environmental forcing incorporates sinusoidal thermal dynamics and persistent spatial autocorrelated noise fields.

\subsection{Emergence}
Community structure, species relative abundances, spatial biomass clustering, elemental dissimilarity indices, and total ecosystem biomass emerge dynamically from individual-level resource uptake, stoichiometric constraints, local competition, and disturbance-driven mortality.

\subsection{Adaptation}
Agents do not possess cognitive or evolutionary adaptive algorithms in the current implementation. Individual elementome $\boldsymbol{q}_i$ are updated dynamically based on ambient elemental uptake ratios and loss processes.

\subsection{Objectives}
Agents do not explicitly solve optimization problems, their growth and reproduction rates are direct programmatic consequences of physiological state, local resource availability, and environmental conditions. They should be maximizing fitness.

\subsection{Learning}
Not implemented.

\subsection{Prediction}
Not implemented.

\subsection{Sensing}
Agents directly sense:
\begin{itemize}
    \item Local inorganic soil nutrient availability $I_{y_i, x_i, k}$.
    \item Local total biomass density $B_{y_i, x_i}$.
    \item Local environmental field values $E_{y_i, x_i}$ and shock field values $S_{y_i, x_i}$.
    \item Local ambient temperature $T_{y_i, x_i}$.
\end{itemize}

\subsection{Interaction}
\begin{itemize}
    \item \textbf{Exploitation Competition:} Agents residing in the same grid cell compete for available inorganic soil nutrients $I_{y,x,k}$.
    \item \textbf{Interference / Space Competition:} Local biomass accumulation $B_{y,x}$ reduces realized growth through a density-dependent space limitation factor $f^{\text{space}} = \max(0, 1 - B_{y,x}/K_{\text{bio}})$.
\end{itemize}

\subsection{Stochasticity}
Stochastic processes are explicitly embedded in:
\begin{enumerate}
    \item Stochastic initialization of inorganic soil pools and agent initial positions.
    \item Drifting nutrient inflow ratios $\boldsymbol{\eta}(t) \sim \mathcal{N}(\boldsymbol{0}, \sigma_{\text{drift}}^2 \mathbf{I})$ and random shock pulses.
    \item Autoregressive spatial Gaussian fields $\mathbf{E}(t)$ and $\mathbf{S}(t)$.
    \item Random spatial location of patch disturbance events $(c_y, c_x)$.
    \item Demographic noise in individual agent biomass updates $\xi_i \sim \mathcal{N}(0, \sigma_{\text{demo}}^2)$.
    \item Probabilistic survival under catastrophe, weak disturbance, and strong disturbance regimes.
    \item Probabilistic seed production $P^{\text{seed}} = 0.1 w_i$, random spatial seed displacement $(\Delta y, \Delta x)$, and probabilistic seedling establishment.
\end{enumerate}

\subsection{Collectives}
Agents of the same species $s$ share species-level traits (niche centers $\boldsymbol{c}_s$, covariance matrices $\boldsymbol{\Sigma}_s$, and seed mass $m_{\text{seed}, s}$). However, collectives do not form explicit computational entities or social structures.

\subsection{Observation}
The model tracks and outputs:
\begin{itemize}
    \item Total active agent count $N_{\text{agents}}(t)$.
    \item Spatial grid of species-specific and total biomass $B_{y,x}(t)$.
    \item Total ecosystem elemental pools $\sum_{i} m_i \boldsymbol{q}_i$.
    \item Mean agent age $\bar{\tau}(t)$ and species-level mean age $\bar{\tau}_s(t)$.
    \item Mean niche fitness of living and dead agents per species.
    \item Per-species nutrient deficit vectors $D_{s, k}(t)$.
    \item Agent-level and species-level elemental dissimilarity indices based on Mahalanobis distance.
\end{itemize}

\section{Details}

\subsection{Initialization}

\subsubsection{Spatial and Soil Initialization}
The soil variable $\mathbf{Soil} \in \mathbb{R}^{H \times W \times 2 N_{\text{soil}}}$ is initialized deterministically and stochastically:
\begin{enumerate}
    \item A base ratio vector $\boldsymbol{r}^{\text{base}} \in \mathbb{R}^{N_{\text{soil}}}$ is normalized such that $\sum_{k=0}^{N_{\text{soil}}-1} r^{\text{base}}_k = 1.0$.
    \item Spatial Gaussian noise $\boldsymbol{\eta}^{\text{soil}}_{y,x,k} \sim \mathcal{N}(0, \sigma_{\text{noise}}^2)$ with $\sigma_{\text{noise}} = 0.05$ is generated.
    \item Spatial noise is smoothed via 2D depthwise convolution with a $3 \times 3$ Gaussian kernel ($\sigma=1.0$), standardized to standard deviation $\sigma_{\text{noise}}$, and blended with raw noise ($\alpha=0.5$).
    \item Soil base ratio with noise is constructed: $\boldsymbol{r}'_{y,x} = \max(0.001, \boldsymbol{r}^{\text{base}} + \mathbf{noise}_{y,x})$.
    \item Normalized soil ratio: $\boldsymbol{r}^{\text{inorg}}_{y,x} = \boldsymbol{r}'_{y,x} / \sum_{k} r'_{y,x,k}$.
    \item Pool size $S_{y,x} = \max(0.5, \mathcal{N}(\mu_{\text{pool}}, \sigma_{\text{pool}}^2))$ with $\mu_{\text{pool}}=1.0, \sigma_{\text{pool}}=0.01$.
    \item Initial inorganic soil pool: $I_{y,x,k}(0) = r^{\text{inorg}}_{y,x,k} \cdot S_{y,x}$.
    \item Initial organic soil pool: $O_{y,x,k}(0) = 0.0$.
\end{enumerate}

\subsubsection{Agent Population Initialization}
Agents are added via \texttt{add\_initial\_seeds(count, species\_id)}:
\begin{itemize}
    \item Coordinates $y_i \sim U(0, H)$, $x_i \sim U(0, W)$.
    \item Species ID $s_i = \text{species\_id}$, initial mass $m_i(0) = 0.1$, age $\tau_i(0) = 0.0$, alive status $\alpha_i(0) = 1.0$.
    \item Elementome vector $\boldsymbol{q}_i(0)$ is generated by sampling raw values $\boldsymbol{q}^{\text{raw}}_i = \boldsymbol{c}_{s_i} \odot \boldsymbol{\xi}$, where $\boldsymbol{\xi} \sim \mathcal{N}(\boldsymbol{1}, 0.05^2 \mathbf{I})$, and normalizing $\boldsymbol{q}_i(0) = \boldsymbol{q}^{\text{raw}}_i / \sum_e q^{\text{raw}}_{i,e}$.
\end{itemize}

\subsection{Input Data}
The model does not rely on external empirical input data files. All inputs, environmental fields, and perturbations are generated dynamically by internal stochastic and deterministic equations.

\subsection{Submodels}

\subsubsection{Soil Biogeochemistry, Spatial Diffusion, and External Inflow}
Inorganic soil nutrients undergo spatial diffusion via a normalized $3 \times 3$ smoothing kernel $K$:
\begin{equation}
K = \begin{bmatrix} 0.05 & 0.10 & 0.05 \\ 0.10 & 0.40 & 0.10 \\ 0.05 & 0.10 & 0.05 \end{bmatrix}
\end{equation}

Symmetrically padded inorganic nutrient pools $\mathbf{I}(t)$ are convolved depthwise:
\begin{equation}
I^{\text{diff}}_{y,x,k}(t) = \sum_{dy=-1}^{1} \sum_{dx=-1}^{1} K_{dy+1, dx+1} \cdot I_{y+dy, x+dx, k}(t)
\end{equation}

External inorganic nutrient inflow is governed by base ratios $\boldsymbol{r}^{\text{base}}$ subjected to drift noise $\boldsymbol{\eta}(t) \sim \mathcal{N}(\boldsymbol{0}, \sigma_{\text{drift}}^2 \mathbf{I})$ with scale $\sigma_{\text{drift}} = 0.08$:
\begin{equation}
r'_k(t) = \max\left(0.001, r^{\text{base}}_k + \eta_k(t)\right), \quad r''_k(t) = \frac{r'_k(t)}{\sum_{j=0}^{N_{\text{soil}}-1} r'_j(t)}
\end{equation}

Pulse events occur with probability $p_{\text{pulse}} = 0.005$. If triggered, pulse indicator $\pi(t)=1$ (else $0$), multiplying inflow by a shock factor $F_{\text{shock}} \sim U(0.1, 3.0)$:
\begin{equation}
r^{\text{in}}_k(t) = r''_k(t) \cdot \left(1 + \pi(t) (F_{\text{shock}} - 1)\right)
\end{equation}
\begin{equation}
I^{\text{new}}_{y,x,k}(t) = I^{\text{diff}}_{y,x,k}(t) + r^{\text{in}}_k(t) \cdot I_{\text{rate}}
\end{equation}
where $I_{\text{rate}} = 0.2$.

Localized patch disturbances occur with probability $p_{\text{dist}} = 0.01$. A center $(c_y, c_x)$ is selected uniformly. For cells where Euclidean distance $d_{y,x} = \sqrt{(y-c_y)^2 + (x-c_x)^2} < R_{\text{dist}}$ ($R_{\text{dist}}=8$), inorganic nutrients are depleted:
\begin{equation}
I^{\text{dist}}_{y,x,k}(t) = I^{\text{new}}_{y,x,k}(t) \cdot \left(1 - S_{\text{dist}} \cdot E_{y,x}(t)\right)
\end{equation}
where $S_{\text{dist}} = 0.7$ and $E_{y,x}(t)$ is the local persistent environmental field. Available soil nutrients are given by:
\begin{equation}
I^{\text{avail}}_{y,x,k}(t) = I^{\text{new}}_{y,x,k}(t) \cdot a_{\text{avail}, k}
\end{equation}

\subsubsection{Environmental Fields and Thermal Dynamics}
Persistent environmental field $\mathbf{E}(t)$ and shock field $\mathbf{S}(t)$ update via an autoregressive spatial process:
\begin{equation}
\mathbf{Z}(t) \sim \mathcal{N}(\boldsymbol{0}, \mathbf{I}), \quad \mathbf{Z}_{\text{smooth}}(t) = \text{AvgPool}_{3\times 3}^{p_{\text{pass}}}\left(\mathbf{Z}(t)\right)
\end{equation}
\begin{equation}
\mathbf{E}'(t) = \rho_{\text{env}} \mathbf{E}(t-1) + (1 - \rho_{\text{env}}) \mathbf{Z}_{\text{smooth}}(t)
\end{equation}
followed by spatial min-max normalisation:
\begin{equation}
E_{y,x}(t) = \frac{E'_{y,x}(t) - \min \mathbf{E}'(t)}{\max \mathbf{E}'(t) - \min \mathbf{E}'(t) + 10^{-8}} \in [0, 1]
\end{equation}

Global temperature $T(t)$ follows a sinusoidal regime:
\begin{equation}
T(t) = T_{\text{mean}} + A_T \sin\left(\frac{2\pi (t + \phi_T)}{P_T}\right)
\end{equation}
where $T_{\text{mean}} = 0.5, A_T = 0.3, P_T = 100.0, \phi_T = 0.0$. Local temperature at grid cell $(y,x)$ is:
\begin{equation}
T_{y,x}(t) = \text{clip}\left(T(t) + \gamma_{\text{sp}} (E_{y,x}(t) - 0.5), 0.0, 1.0\right)
\end{equation}

\subsubsection{Niche Fitness and Trait-Based Physiology}
For agent $i$ with elementome $\boldsymbol{q}_i(t)$ and species identity $s_i$, niche fitness $w_i(t)$ is computed using the Mahalanobis distance to species niche center $\boldsymbol{c}_{s_i}$:
\begin{equation}
\boldsymbol{\delta}_i(t) = \boldsymbol{q}_i(t) - \boldsymbol{c}_{s_i}
\end{equation}
\begin{equation}
d_{M, i}^2(t) = \boldsymbol{\delta}_i(t)^T \boldsymbol{\Sigma}_{s_i}^{-1} \boldsymbol{\delta}_i(t)
\end{equation}
\begin{equation}
w_i(t) = \text{clip}\left(1.0 - \left(\frac{\sqrt{d_{M, i}^2(t)}}{\sigma_{\text{thresh}}}\right)^2, 0.0, 1.0\right)
\end{equation}
where $\sigma_{\text{thresh}} = 3.0$.

\subsubsection{Nutrient Uptake, Quota Bottlenecks, and Growth}
Desired agent growth depends on niche fitness, intrinsic growth rate $g_0 = 0.7$, biomass $m_i(t)$, and local temperature growth factor $F_g(T_i) = 1 + \beta_g (T_i - 0.5)$ ($\beta_g = 0.5$):
\begin{equation}
G_i^{\text{des}}(t) = w_i(t) \cdot g_0 \cdot m_i(t) \cdot \left(1 + \beta_g (T_{y_i, x_i}(t) - 0.5)\right)
\end{equation}

Carbon uptake is self-supplied via internal carbon quota fraction $q_{i,0}(t)$:
\begin{equation}
U_{i, \text{C}}^{\text{pot}}(t) = G_i^{\text{des}}(t) \cdot q_{i,0}(t)
\end{equation}
Remaining growth demand $R_i(t) = G_i^{\text{des}}(t) - U_{i, \text{C}}^{\text{pot}}(t)$ is allocated across soil elements according to normalized species soil niche preferences $\tilde{\boldsymbol{n}}_{s_i}$:
\begin{equation}
\tilde{n}_{s_i, k} = \frac{c_{s_i, k+1}}{\sum_{j=1}^{N_{\text{soil}}} c_{s_i, j} + 10^{-9}}, \quad k \in \{0, \dots, N_{\text{soil}}-1\}
\end{equation}
\begin{equation}
U_{i, k}^{\text{des}}(t) = R_i(t) \cdot \tilde{n}_{s_i, k}
\end{equation}

Soil nutrient uptake follows Michaelis-Menten kinetics with half-saturation constant $K_m = 0.1$:
\begin{equation}
U_{i, k}(t) = U_{i, k}^{\text{des}}(t) \cdot \frac{I^{\text{avail}}_{y_i, x_i, k}(t)}{I^{\text{avail}}_{y_i, x_i, k}(t) + K_m}
\end{equation}

Elemental mass pools after uptake are:
\begin{equation}
\mathbf{E}_i^{\text{pool}}(t) = m_i(t) \boldsymbol{q}_i(t) + \mathbf{U}_i^{\text{full}}(t)
\end{equation}
where $U_{i, 0}^{\text{full}} = U_{i, \text{C}}^{\text{pot}}$ and $U_{i, k+1}^{\text{full}} = U_{i, k}$.

Cell quotas are $Q_{i, e}(t) = E_{i, e}^{\text{pool}}(t) / (m_i(t) + 10^{-9})$. Growth limitation obeys Liebig's Law of the Minimum relative to minimum quota $Q_{\min} = 0.01$:
\begin{equation}
g_{i, e}(t) = \max\left(0, 1 - \frac{Q_{\min}}{Q_{i, e}(t) + 10^{-9}}\right), \quad g_i^{\text{lim}}(t) = \min_{e \in \{0,\dots,N_{\text{org}}-1\}} g_{i, e}(t)
\end{equation}

Density-dependent space limitation factor depends on total local biomass $B_{y_i, x_i}(t) = \sum_{j: (y_j, x_j)=(y_i, x_i)} m_j(t)$ relative to carrying capacity $K_{\text{bio}} = 1.5$:
\begin{equation}
f_i^{\text{space}}(t) = \max\left(0, 1 - \frac{B_{y_i, x_i}(t)}{K_{\text{bio}}}\right)
\end{equation}

Realized growth is:
\begin{equation}
G_i^{\text{real}}(t) = \left(U_{i, \text{C}}(t) + \sum_{k=0}^{N_{\text{soil}}-1} U_{i, k}(t)\right) \cdot g_i^{\text{lim}}(t) \cdot f_i^{\text{space}}(t)
\end{equation}

Maintenance respiration uses temperature factor $F_r(T_i) = 1 + \beta_r (T_i - 0.5)$ ($\beta_r = 0.3$) and rate $r_{\text{resp}} = 0.02$:
\begin{equation}
M_i(t) = m_i(t) \cdot r_{\text{resp}} \cdot \left(1 + \beta_r (T_{y_i, x_i}(t) - 0.5)\right)
\end{equation}

Biomass updates with additive demographic noise $\xi_i(t) \sim \mathcal{N}(0, \sigma_{\text{demo}}^2)$ ($\sigma_{\text{demo}} = 0.003$):
\begin{equation}
m_i''(t) = m_i(t) + G_i^{\text{real}}(t) - M_i(t) + \xi_i(t)
\end{equation}
Initial viability requires $m_i''(t) > 0.01$, yielding initial alive flag $\alpha_i^{(0)}(t) = \mathbb{I}(m_i''(t) > 0.01)$.

\subsubsection{Mortality and Multi-Tier Disturbance Regimes}
Surviving alive status $\alpha_i(t)$ is sequentially filtered through three disturbance mechanisms:
\begin{enumerate}
    \item \textbf{Catastrophe Event} (if $t \pmod{T_{\text{cat}}} == 0$):
    \begin{equation}
    P_i^{\text{cat}}(t) = d_{\text{cat}} \cdot S_{y_i, x_i}(t), \quad \alpha_i^{(1)}(t) = \alpha_i^{(0)}(t) \cdot \mathbb{I}\left(u_{i, \text{cat}} > P_i^{\text{cat}}(t)\right)
    \end{equation}
    where $d_{\text{cat}} = 0.4$, $T_{\text{cat}} = 200$, and $u_{i,\text{cat}} \sim U(0,1)$.
    
    \item \textbf{Weak Disturbance Event} (if $T_{\text{weak}} > 0$ and $t \pmod{T_{\text{weak}}} == 0$):
    \begin{equation}
    P_i^{\text{weak}}(t) = d_{\text{weak}} \cdot E_{y_i, x_i}(t) \cdot \left(1 - w_i(t)\right)^{\alpha_{\text{weak}}}
    \end{equation}
    \begin{equation}
    \alpha_i^{(2)}(t) = \alpha_i^{(1)}(t) \cdot \mathbb{I}\left(u_{i, \text{weak}} > P_i^{\text{weak}}(t)\right)
    \end{equation}
    where $d_{\text{weak}} = 0.4, \alpha_{\text{weak}} = 2.0$.
    
    \item \textbf{Strong Disturbance Event} (if $T_{\text{strong}} > 0$ and $t \pmod{T_{\text{strong}}} == 0$):
    \begin{equation}
    \theta_i(t) = \frac{m_i''(t) - \min_j m_j''(t)}{\max_j m_j''(t) - \min_j m_j''(t) + 10^{-9}}
    \end{equation}
    \begin{equation}
    P_i^{\text{strong}}(t) = d_{\text{strong}} \cdot S_{y_i, x_i}(t) \cdot \left(\theta_i(t)\right)^{\beta_{\text{strong}}}
    \end{equation}
    \begin{equation}
    \alpha_i(t) = \alpha_i^{(2)}(t) \cdot \mathbb{I}\left(u_{i, \text{strong}} > P_i^{\text{strong}}(t)\right)
    \end{equation}
    where $d_{\text{strong}} = 0.4, \beta_{\text{strong}} = 2.0$.
\end{enumerate}

\subsubsection{Turnover, Organic Decay, and Mineralization}
Surviving agents undergo turnover loss at rate $\mu_{\text{turn}} = 0.02$. Total elemental mass lost per agent from mortality and turnover is:
\begin{equation}
L_{i, e}(t) = \left(m_i^+(t) (1 - \alpha_i(t)) + m_i^+(t) \mu_{\text{turn}} \alpha_i(t)\right) \cdot Q_{i, e}(t)
\end{equation}
where $m_i^+(t) = \max(0, m_i''(t))$.

Surviving agent final biomass and normalized elemental ratios update as:
\begin{equation}
m_i(t+1) = \left(m_i''(t) - m_i^+(t) \mu_{\text{turn}}\right) \cdot \alpha_i(t)
\end{equation}
\begin{equation}
q'_{i, e}(t+1) = \text{clip}\left(\frac{E_{i, e}^{\text{pool}}(t) - L_{i, e}(t)}{m_i(t+1) + 10^{-9}}, 0.01, 0.99\right), \quad q_{i, e}(t+1) = \frac{q'_{i, e}(t+1)}{\sum_j q'_{i, j}(t+1)}
\end{equation}

Soil recycled organic inputs per grid cell are $F^{\text{fresh}}_{y,x,k}(t) = \sum_{i: (y_i,x_i)=(y,x)} L_{i, k+1}(t)$. Total organic soil pool updates with temperature-dependent mineralization rate $m_{\text{min}} = 0.05$:
\begin{equation}
O^{\text{tot}}_{y,x,k}(t) = O_{y,x,k}(t) + F^{\text{fresh}}_{y,x,k}(t)
\end{equation}
\begin{equation}
\Phi^{\text{min}}_{y,x,k}(t) = O^{\text{tot}}_{y,x,k}(t) \cdot m_{\text{min}} \cdot \left(1 + \beta_m (T_{y,x}(t) - 0.5)\right)
\end{equation}
\begin{equation}
O_{y,x,k}(t+1) = O^{\text{tot}}_{y,x,k}(t) - \Phi^{\text{min}}_{y,x,k}(t)
\end{equation}
\begin{equation}
I_{y,x,k}(t+1) = \max\left(0, I^{\text{new}}_{y,x,k}(t) + \Phi^{\text{min}}_{y,x,k}(t) - \sum_{i: (y_i,x_i)=(y,x)} U_{i, k}(t)\right)
\end{equation}

\subsubsection{Reproduction, Dispersal, and Establishment}
Surviving agents with biomass $m_i(t+1) > C_{\text{seed}}$ ($C_{\text{seed}} = 0.3$) are eligible to reproduce with probability $P_i^{\text{seed}}(t) = 0.1 w_i(t)$. If reproducing, parent biomass is reduced: $m_i(t+1) \leftarrow m_i(t+1) - C_{\text{seed}}$.

Seed dispersal distance range per species $s_i$ is:
\begin{equation}
R_{s_i} = R_{\text{scale}} \cdot \left(m_{\text{seed}, s_i} + \epsilon_{\text{rng}}\right)^{\alpha_{\text{rng}}}
\end{equation}
where $R_{\text{scale}} = 10.0, \alpha_{\text{rng}} = 1.0, \epsilon_{\text{rng}} = 10^{-6}$.

Grid offset radius $d_{\text{spread}} = \max(1, \lfloor \text{round}(R_{s_i}) \rfloor)$. Spatial integer offsets $\Delta y, \Delta x$ are drawn uniformly from $\{-d_{\text{spread}}, \dots, d_{\text{spread}}\}$. Offspring coordinates on the toroidal grid are:
\begin{equation}
y_{\text{off}} = (y_i + \Delta y) \pmod H, \quad x_{\text{off}} = (x_i + \Delta x) \pmod W
\end{equation}

Seedling establishment probability depends on pre-existing local biomass density $B_{y_{\text{off}}, x_{\text{off}}}$:
\begin{equation}
P^{\text{est}} = \max\left(0, 1 - \frac{B_{y_{\text{off}}, x_{\text{off}}}}{K_{\text{bio}}}\right)
\end{equation}
Established seedlings inherit the parent's updated elementome ratio $\boldsymbol{q}_i(t+1)$ and start with species seed mass $m_{\text{seed}, s_i}$ and age $\tau = 0$.

\section{Mathematical Analysis and Conservation Properties}

\subsection{Elemental Conservation Analysis}
To evaluate whether elemental mass is strictly conserved in the model, consider the total mass $M_k^{\text{tot}}(t)$ of soil element $k \in \{0, \dots, N_{\text{soil}}-1\}$ across the entire system:
\begin{equation}
M_k^{\text{tot}}(t) = \sum_{y=0}^{H-1} \sum_{x=0}^{W-1} \left( I_{y,x,k}(t) + O_{y,x,k}(t) \right) + \sum_{i=1}^{M} \alpha_i(t) m_i(t) q_{i, k+1}(t)
\end{equation}

Summing the explicit update equations yields:
\begin{equation}
M_k^{\text{tot}}(t+1) - M_k^{\text{tot}}(t) = \sum_{y,x} \left( r^{\text{in}}_k(t) I_{\text{rate}} - \text{Disturbance Losses} \right) - \sum_{i \in \text{Repro}} C_{\text{seed}} q_{i, k+1} + \sum_{\text{offspring}} m_{\text{seed}, s_i} q_{i, k+1}
\end{equation}

\begin{quote}
\textbf{Mathematical Implication on Conservation Laws:} The model does \textbf{NOT} satisfy global mass conservation. Elemental mass enters non-conservatively via external atmospheric inflow ($I_{\text{rate}}$), random pulse shocks, and seed creation when $C_{\text{seed}} \neq m_{\text{seed}, s_i}$. Conversely, elemental mass is lost via localized patch disturbances and clipping constraints on soil pools ($I_{y,x,k} \ge 0$).
\end{quote}

\section{Parameters and State Variables}

\begin{longtable}{lllp{4.5cm}l}
\caption{Complete Specification of Model Parameters.} \label{tab:parameters} \\
\toprule
\textbf{Symbol} & \textbf{Code Parameter} & \textbf{Default Value} & \textbf{Meaning / Role} & \textbf{Units} \\
\midrule
\endhead
$H$ & \texttt{height} & -- & Grid height & cells \\
$W$ & \texttt{width} & -- & Grid width & cells \\
$M$ & \texttt{max\_agents} & -- & Maximum agent buffer size & agents \\
$g_0$ & \texttt{growth\_rate} & $0.7$ & Base growth rate & $\text{step}^{-1}$ \\
$r_{\text{resp}}$ & \texttt{respiration\_rate} & $0.02$ & Base maintenance respiration rate & $\text{step}^{-1}$ \\
$\mu_{\text{turn}}$ & \texttt{turnover\_rate} & $0.02$ & Biomass turnover fraction & $\text{step}^{-1}$ \\
$m_{\text{min}}$ & \texttt{mineralization\_rate} & $0.05$ & Base organic mineralization rate & $\text{step}^{-1}$ \\
$C_{\text{seed}}$ & \texttt{seed\_cost} & $0.3$ & Biomass cost to produce a seed & mass \\
$m_{\text{seed}}$ & \texttt{seed\_mass} & $0.05$ & Default seed biomass & mass \\
$K_{\text{bio}}$ & \texttt{K\_biomass} & $1.5$ & Local biomass carrying capacity & mass $\text{cell}^{-1}$ \\
$I_{\text{rate}}$ & \texttt{soil\_input\_rate} & $0.2$ & Inorganic nutrient inflow rate & mass $\text{cell}^{-1} \text{step}^{-1}$ \\
$\sigma_{\text{drift}}$ & \texttt{input\_drift\_scale} & $0.08$ & Inflow ratio drift std dev & dimensionless \\
$\sigma_{\text{thresh}}$ & \texttt{sigma\_threshold} & $3.0$ & Mahalanobis niche fitness threshold & dimensionless \\
$T_{\text{cat}}$ & \texttt{catastrophe\_interval} & $200$ & Catastrophe event interval & timesteps \\
$d_{\text{cat}}$ & \texttt{catastrophe\_mortality} & $0.4$ & Catastrophe mortality multiplier & dimensionless \\
$p_{\text{dist}}$ & \texttt{p\_disturbance} & $0.01$ & Patch disturbance probability & $\text{step}^{-1}$ \\
$R_{\text{dist}}$ & \texttt{disturbance\_radius} & $8$ & Patch disturbance radius & cells \\
$S_{\text{dist}}$ & \texttt{disturbance\_strength} & $0.7$ & Soil depletion fraction & dimensionless \\
$\sigma_{\text{demo}}$ & \texttt{demo\_noise\_std} & $0.003$ & Demographic biomass noise std dev & mass \\
$\rho_{\text{env}}$ & \texttt{env\_field\_persistence} & $0.85$ & Environmental field persistence & dimensionless \\
$\rho_{\text{shock}}$ & \texttt{shock\_field\_persistence} & $0.85$ & Shock field persistence & dimensionless \\
$T_{\text{mean}}$ & \texttt{temperature\_mean} & $0.5$ & Mean global temperature & temperature units \\
$A_T$ & \texttt{temperature\_amplitude} & $0.3$ & Temperature sine amplitude & temperature units \\
$P_T$ & \texttt{temperature\_period} & $100.0$ & Temperature period & timesteps \\
$\beta_g$ & \texttt{temperature\_growth\_strength} & $0.5$ & Growth temperature sensitivity & dimensionless \\
$\beta_r$ & \texttt{temperature\_respiration\_strength} & $0.3$ & Respiration temperature sensitivity & dimensionless \\
$\beta_m$ & \texttt{temperature\_mineralization\_strength} & $0.4$ & Mineralization temp sensitivity & dimensionless \\
\bottomrule
\end{longtable}

\begin{longtable}{lllp{4cm}lp{2.5cm}}
\caption{Complete Specification of Primary State Variables.} \label{tab:state_variables} \\
\toprule
\textbf{Symbol} & \textbf{Code Variable} & \textbf{Type} & \textbf{Description} & \textbf{Units} & \textbf{Domain} \\
\midrule
\endhead
$y_i$ & \texttt{agents[:, 0]} & Agent & $y$-coordinate & cells & $[0, H)$ \\
$x_i$ & \texttt{agents[:, 1]} & Agent & $x$-coordinate & cells & $[0, W)$ \\
$s_i$ & \texttt{agents[:, 2]} & Agent & Species identity index & ID & $\{0, \dots, N_{\text{spp}}-1\}$ \\
$m_i$ & \texttt{agents[:, 3]} & Agent & Individual biomass & mass & $[0, \infty)$ \\
$\boldsymbol{q}_i$ & \texttt{agents[:, 4:15]} & Agent & Elementome ($N_{\text{org}}$) & fraction & $[0, 1]^{N_{\text{org}}}$ \\
$\alpha_i$ & \texttt{agents[:, 15]} & Agent & Viability status indicator & binary & $\{0.0, 1.0\}$ \\
$\tau_i$ & \texttt{agents[:, 16]} & Agent & Individual agent age & timesteps & $[0, \infty)$ \\
$I_{y,x,k}$ & \texttt{soil[:, :, 0:10]} & Soil Grid & Inorganic nutrient pool & mass $\text{cell}^{-1}$ & $[0, \infty)$ \\
$O_{y,x,k}$ & \texttt{soil[:, :, 10:20]} & Soil Grid & Organic nutrient pool & mass $\text{cell}^{-1}$ & $[0, \infty)$ \\
$E_{y,x}$ & \texttt{env\_field} & Env Field & Persistent environmental noise & dimensionless & $[0, 1]$ \\
$S_{y,x}$ & \texttt{shock\_field} & Env Field & Persistent shock noise & dimensionless & $[0, 1]$ \\
$T$ & \texttt{temperature} & Global Env & Global thermal state & dimensionless & $[0, 1]$ \\
\bottomrule
\end{longtable}

\section{Outputs and Observations}
The model computes and logs summary observation metrics:
\begin{enumerate}
    \item \textbf{Species Biomass Grid:} $B_{y,x,s}(t) = \sum_{i: (y_i,x_i)=(y,x), s_i=s} \alpha_i m_i$.
    \item \textbf{Total Ecosystem Elemental Pools:} $P_e(t) = \sum_{i: \alpha_i=1} m_i q_{i, e}$.
    \item \textbf{Elemental Dissimilarity Indices:} Mahalanobis distance-based diversity metrics computed across agents (\texttt{get\_agent\_elemental\_dissimilarity\_index\_tf}) and across species weighted by biomass (\texttt{get\_species\_elemental\_dissimilarity\_index\_tf}):
    \begin{equation}
    \text{EDm} = \frac{1}{E} \sum_{i, j} D_M(\boldsymbol{q}_i, \boldsymbol{q}_j) \cdot \min(p_i, p_j)
    \end{equation}
\end{enumerate}

\section{Reproducibility}
\begin{itemize}
    \item \textbf{Programming Language \& Libraries:} Python 3.x, TensorFlow 2.x, NumPy.
    \item \textbf{Execution Paradigm:} Dynamic TensorFlow graph execution via \texttt{@tf.function} vectorized tensor operations.
    \item \textbf{Seed Handling:} Pseudorandom generation relies on TensorFlow/NumPy global state; explicit random seed initialization is recommended prior to simulation runs.
\end{itemize}

\section{Code-to-Mathematics Traceability}

\begin{longtable}{lp{4.5cm}p{5.5cm}l}
\caption{Code-to-Mathematics Traceability Matrix.} \label{tab:traceability} \\
\toprule
\textbf{Model Component} & \textbf{Mathematical Symbol / Equation} & \textbf{TensorFlow Code Location / Method} & \textbf{Line Reference} \\
\midrule
\endhead
Soil Diffusion & $I^{\text{diff}}_{y,x,k} = \text{Conv2D}(I, K)$ & \texttt{tf.nn.depthwise\_conv2d(inorg\_padded, ...)} & \texttt{step()} \\
Drifted Inflow & $\boldsymbol{r}' = \max(0.001, \boldsymbol{r}^{\text{base}} + \boldsymbol{\eta})$ & \texttt{input\_ratio\_drifted = tf.maximum(...)} & \texttt{step()} \\
Environmental Fields & $\mathbf{E}' = \rho \mathbf{E} + (1-\rho) \mathbf{Z}_{\text{smooth}}$ & \texttt{\_update\_spatial\_field()} & Method Definition \\
Niche Fitness & $w_i = 1 - (d_M / \sigma_{\text{thresh}})^2$ & \texttt{\_compute\_niche\_fitness\_mahalanobis()} & Method Definition \\
Michaelis-Menten Uptake & $U_{i,k} = U^{\text{des}}_{i,k} I / (I + K_m)$ & \texttt{uptake = desired\_soil * (available / (available + K\_m))} & \texttt{step()} \\
Liebig Quota Constraint & $g_{i,e} = \max(0, 1 - Q_{\min}/Q_{i,e})$ & \texttt{g = tf.maximum(0.0, 1.0 - (MIN\_QUOTA / Q))} & \texttt{step()} \\
Space Constraint & $f^{\text{space}} = \max(0, 1 - B / K_{\text{bio}})$ & \texttt{space\_f = tf.maximum(0.0, 1.0 - (loc\_bio / K\_biomass))} & \texttt{step()} \\
Respiration & $M_i = m_i r_{\text{resp}} (1 + \beta_r (T_i - 0.5))$ & \texttt{maint = mass * respiration\_rate * temp\_resp\_factor} & \texttt{step()} \\
Mineralization Flux & $\Phi^{\text{min}} = O^{\text{tot}} m_{\text{min}} F_m(T)$ & \texttt{flux = org\_tot * mineralization\_rate * temp\_mineral\_factor} & \texttt{step()} \\
Dispersal Torus & $y_{\text{off}} = (y_i + \Delta y) \pmod H$ & \texttt{ny = (tf.cast(p\_dat[:, I\_Y], int32) + dy) \% self.H} & \texttt{step()} \\
\bottomrule
\end{longtable}

\subsection{Detailed Model Architecture and Visual Execution Workflow}

The structural and dynamical organization of the \texttt{HybridEcosystem} simulation is formalised in Figure~\ref{fig:model_architecture}. The model operates as a multi-tier complex adaptive system that couples abiotic forcing, 2D spatial diffusion, individual-level elemental physiological updates, multi-tier disturbance regimes, and ecosystem-level elemental recycling.

\begin{figure}[H]
\centering
% Insert TikZ Diagram Here
\caption{\textbf{Comprehensive ODD 2020 Architectural Model Diagram of \texttt{HybridEcosystem}.} The diagram illustrates the five structural levels: Level 1 (External Forcing and Parameter Inputs), Level 2 (Abiotic Soil Grid and Spatial Fields), Level 3 (Individual Agent Buffer and State Vector), Level 4 (Sequential Process Submodels with governing equations [P1]--[P10]), and Level 5 (Emergent System States and Logged Outputs). Arrow styles specify causality, exogenous forcing, stochastic draws, state transitions, and closed feedback loops.}
\label{fig:model_architecture}
\end{figure}

\subsubsection{Structural Hierarchy and Information Flow}
The simulation execution within each discrete step $t \to t+1$ follows a strict hierarchical information flow across five structural levels:

\begin{enumerate}
    \item \textbf{Level 1 (Exogenous Forcing \& Traits):} External parameters, species stoichiometric niche centers $\boldsymbol{c}_s$, covariance matrices $\boldsymbol{\Sigma}_s$, sinusoidal global thermal regimes $T(t)$, and stochastic atmospheric drift vectors $\boldsymbol{\eta}(t)$ drive the simulation without feedback from internal system state.
    
    \item \textbf{Level 2 (Abiotic Soil Grid):} The abiotic state consists of a regular 2D lattice ($H \times W$) containing $N_{\text{soil}}$ inorganic pools $\mathbf{I}(t)$ and organic pools $\mathbf{O}(t)$, along with spatial fields for environmental noise $E_{y,x}$ and physical shocks $S_{y,x}$. Inorganic nutrients undergo spatial convolution diffusion:
    \begin{equation}
    I^{\text{diff}}_{y,x,k}(t) = \sum_{dy=-1}^{1} \sum_{dx=-1}^{1} K_{dy+1, dx+1} \cdot I_{y+dy, x+dx, k}(t)
    \end{equation}
    where $K$ is a normalized $3 \times 3$ Gaussian smoothing kernel.
    
    \item \textbf{Level 3 (Individual Agent Buffer):} The biotic component is represented by an agent buffer $\mathbf{A}(t)$ holding state vectors $\mathbf{a}_i(t) = (y_i, x_i, s_i, m_i, \boldsymbol{q}_i, \alpha_i, \tau_i)^T$ for every active organism $i$.
    
    \item \textbf{Level 4 (Sequential Submodels [P1]--[P10]):} State transitions occur via ten sequential algorithmic blocks:
    \begin{itemize}
        \item \textbf{[P1] Soil Diffusion and Inflow:} Updates inorganic soil pools via depthwise 2D spatial convolution and drifted atmospheric input.
        \item \textbf{[P2] Environmental Field Autoregression:} Updates persistent noise field $E_{y,x}$ and shock field $S_{y,x}$ via spatial autoregressive smoothing, and evaluates local cell temperature $T_{y,x}(t) = \text{clip}(T(t) + \gamma_{\text{sp}}(E_{y,x} - 0.5), 0, 1)$.
        \item \textbf{[P3] Localized Patch Disturbance:} Depletes inorganic nutrients within radius $R_{\text{dist}}$ around stochastic centers $(c_y, c_x)$ proportionally to $S_{\text{dist}} E_{y,x}$.
        \item \textbf{[P4] Niche Fitness Evaluation:} Computes Mahalanobis distance $d_{M,i}$ between elementome $\boldsymbol{q}_i$ and species niche optimum $\boldsymbol{c}_{s_i}$:
        \begin{equation}
        w_i = \text{clip}\left(1.0 - \frac{(\boldsymbol{q}_i - \boldsymbol{c}_{s_i})^T \boldsymbol{\Sigma}_{s_i}^{-1} (\boldsymbol{q}_i - \boldsymbol{c}_{s_i})}{\sigma_{\text{thresh}}^2}, 0.0, 1.0\right)
        \end{equation}
        \item \textbf{[P5] Michaelis-Menten Resource Uptake:} Calculates nutrient uptake $U_{i,k}$ based on saturation kinetics with half-saturation constant $K_m = 0.1$:
        \begin{equation}
        U_{i,k} = U^{\text{des}}_{i,k} \cdot \frac{I^{\text{avail}}_{y_i, x_i, k}}{I^{\text{avail}}_{y_i, x_i, k} + K_m}
        \end{equation}
        \item \textbf{[P6] Quota and Space Limitation:} Bottlenecks growth using Liebig's Law of the Minimum across internal element quotas $Q_{i,e}$ ($g_i^{\text{lim}} = \min_e \max(0, 1 - Q_{\min}/Q_{i,e})$) and density-dependent spatial crowding ($f_i^{\text{space}} = \max(0, 1 - B_{y_i,x_i}/K_{\text{bio}})$).
        \item \textbf{[P7] Realized Growth and Respiration:} Updates individual biomass $m_i'' = m_i + G_i^{\text{real}} - M_i + \xi_i$ with temperature-dependent respiration $M_i$ and additive demographic noise $\xi_i \sim \mathcal{N}(0, \sigma_{\text{demo}}^2)$.
        \item \textbf{[P8] Multi-Tier Mortality Filtering:} Filters agent survival through sequential viability, periodic catastrophe, weak niche-dependent, and strong mass-dependent mortality filters.
        \item \textbf{[P9] Turnover Recycling and Mineralization:} Recycles elemental losses $L_{i,e}$ into organic soil pools $O_{y,x,k}$, which undergo temperature-dependent mineralization to replenish inorganic pools:
        \begin{equation}
        \Phi^{\text{min}}_{y,x,k} = O^{\text{tot}}_{y,x,k} \cdot m_{\text{min}} \cdot \left(1 + \beta_m (T_{y,x} - 0.5)\right)
        \end{equation}
        \item \textbf{[P10] Dispersal and Establishment:} Disperses offspring seeds on a 2D torus ($y_{\text{off}} = (y_i + \Delta y) \pmod H, x_{\text{off}} = (x_i + \Delta x) \pmod W$) subject to density-dependent establishment probability $P^{\text{est}} = \max(0, 1 - B_{\text{off}}/K_{\text{bio}})$.
    \end{itemize}
    
    \item \textbf{Level 5 (Emergent System Outputs):} Computes aggregate community metrics including species spatial biomass grids $B_{y,x,s}(t)$, global elemental pools $P_e(t)$, and Mahalanobis elemental dissimilarity index $\text{EDm}$.
\end{enumerate}

\subsubsection{Causal Feedback Dynamics}
The architectural diagram explicitly highlights three primary closed feedback loops that govern long-term community dynamics:

\begin{enumerate}
    \item \textbf{Resource Depletion Loop (Negative Feedback):} Active uptake by dense agent clusters reduces local inorganic nutrient availability $I^{\text{avail}}_{y,x,k}$, which suppresses subsequent uptake $U_{i,k}$ via Michaelis-Menten saturation, stabilizing local biomass accumulation.
    
    \item \textbf{Nutrient Recycling Loop (Positive/Balancing Feedback):} Agent turnover and mortality transfer elemental mass to organic pools $O_{y,x,k}$. Thermal mineralization converts this organic matter back into inorganic nutrients $I_{y,x,k}$, creating a localized nutrient feedback loop that sustains productivity in established patches.
    
    \item \textbf{Biomass Crowding Loop (Negative Feedback):} High local biomass density $B_{y,x}$ directly restricts individual growth via $f_i^{\text{space}}$ and reduces seedling establishment success via $P^{\text{est}}$, preventing unbounded local population growth and promoting spatial self-organization.
\end{enumerate}



\end{document}